\section{Related Work}

Beyond the general summarization tasks and methods covered in~\Cref{sec:background-tasks,sec:background-methods}, there are several areas of related research relevant to the problems of templatic document transformation and evaluation.

Although the tasks of slide and poster generation have generally been considered separate from scientific summarization, they are related in that both tasks require taking an input article, then organizing and abstracting the information to generate a new document. Past work has developed methods for slide generation from papers~\cite{Hu2015PPSGenLP, Li2021TowardsTS, Hu2015PPSGenLP,Fu2021DOC2PPTAP}, from code~\cite{Wang2023Slide4NCP}, or based on a query~\cite{Sun2021D2SDG}. Poster generation has been explored in the form of content extraction for posters~\cite{Xu2021NeuralCE}, interactive generation~\cite{Wang2015iPosterIP}, or full content generation using graphical models~\cite{Qiang2016LearningTG}. To the best of our knowledge, our work is the first to create a unified method capable of generating a diverse range of templatic views of a source document.

Similar to chain of thought prompting~\cite{Wei2022ChainOT} and content planning prompting~\cite{Wang2022SelfConsistencyIC}, we show that generating an intermediate representation of an input document can improve performance over simply prompting the model to generate the final document from the original input.

As past work has tackled generation of templatic views as separate tasks, methods for automatic evaluation of different document types is fragmented. LongSumm, the shared task introduced by~\textcite{chandrasekaran2020overview}, uses ROUGE to evaluate model performance~\cite{lin2004rouge}. \textcite{Fu2021DOC2PPTAP} introduced Slide Level ROUGE to evaluate slide generation, a variant that contains a penalty for the number of slides. \textcite{Qiang2016LearningTG} used a trained regressor. The general summarization metrics surveyed in~\Cref{sec:background-eval-summaries} are intended for a simple input document-summary setup, and do not take into account factors that affect the quality of other types of documents (e.g. structure). Our work is the first to introduce template adaptable evaluation, allowing uniform comparison of performance across template types.
